% LSPV-1 REPORT::TREE-FACTOR-JOIN135::2026-10-25T02:30:00+09:00
\documentclass[11pt]{article}
\usepackage[T1]{fontenc}
\usepackage[margin=.85in]{geometry}
\usepackage{amsmath,amssymb,amsthm,booktabs,longtable,xurl,hyperref}
\newtheorem{theorem}{Theorem}
\newtheorem{lemma}[theorem]{Lemma}
\newcommand{\SAT}{\operatorname{SAT}}
\title{TREE-FACTOR-JOIN135: A Tree of Signed Interval-Factor Cores\\
and the Precise Remaining All-Input ARC Bridge}
\author{LSPV-1: Spine Lead and Archive Reconciliation}
\date{25 October 2026, 02:30 Asia/Seoul}
\begin{document}\maketitle
\begin{center}
\textbf{P versus NP remains OPEN. ARC-EDGE-90 remains OPEN.}\\
\textbf{PROVED/EXACT within a certified restricted source-decomposition model,}
independent review pending. \textbf{COMPUTATIONAL} finite controls only.
\end{center}

\section{Target, precise endpoint and three falsifiable routes}
Let $\mathcal E$ be the set of simultaneous exponent pairs $(t,s)$ of one
fixed uniform exact ordinary-SAT decider, correct on every encoded input
of bit length $N$, in the indexed read-only-input random-access Turing
machine model with counted writable/index tapes. Under $P=NP$, some finite
pair lies in $\mathcal E$. It would suffice to prove a cofinal finite
charged route of achievable pairs from arbitrarily large diagonal
$(K,K)\in\mathcal E$ to a genuinely forbidden endpoint.
Corollary 5.8 of Fortnow--Lipton--van Melkebeek--Viglas (FLMV) gives
$\operatorname{NTIME}[n]\not\subseteq
\operatorname{DTISP}[n^c,o(n^{\beta(c)})]$, where
\[
\beta(c)=\tfrac12((c+2)/c^2-c),\qquad
c_\star=241/200,\quad
\beta(c_\star)=11642479/23232400>1/2.
\]
A fixed all-input joint exponent pair with $t<c_\star$ and
$s<\beta(c_\star)$ would be forbidden, with strict slack for the SAT
transfer. The old universal $t+s<2\cos(\pi/7)$ landing is
\textbf{FALSIFIED/RETIRED}. No favorable all-input edge has been proved.

Three mechanisms were specified before testing:
\begin{description}
\item[M1.] Combine Role 1's signed degree-two interval-factor solver with
Role 5's exact named-frontier gluing on a tree of blocks rather than a
single scoped path. \textbf{PROVED/EXACT restricted}, below.
\item[M2.] Permit a \emph{source-defined} number $k_v$ of internal
occurrence-degree-greater-than-two variables per block, charging their
assignments. \textbf{PROVED/EXACT restricted}, below.
\item[M3.] Infer a favorable ARC successor, or a mandatory decision-time
lower bound, merely from the boundary count or the factor table size.
\textbf{NOT ESTABLISHED}; it is unsound as a universal inference because
an exact solver need not use the declared decomposition and easy affine
and Horn sources can have arbitrarily large interfaces.
\end{description}

\section{Typed source and a certified tree of blocks}
An instance $H$ is a finite list of \emph{named} signed interval
cardinality constraints $(C_j,L_j,U_j)$, with pairwise distinct variables
within each constraint, demanding
$L_j\le\sum_{\ell\in C_j}[\ell=1]\le U_j$.
Ordinary CNF is included by $(L_j,U_j)=(1,|C_j|)$, and signed exact-one
by $(1,1)$. A source certificate partitions constraint indices among
blocks $V(T)$ and supplies an undirected tree $T$ on the blocks.
For each variable $x$, the set
$I_x=\{v:x\text{ occurs in a constraint assigned to }v\}$ must
be a connected subtree. This is the \emph{running-intersection contract}.
It is explicitly checked, not an assumed consequence of clause count.

Write $X_v$ for source names occurring in block $v$,
$B_v=\{x\in X_v:x\in X_u\text{ for some }u\ne v\}$,
and $S_{vu}=X_v\cap X_u$ for adjacent blocks.
For each $v$, let $k_v$ count the \emph{internal} names
$x\in X_v\setminus B_v$ occurring in more than two constraints of block
$v$. These are source-defined measures, independent of runtime and the
SAT/UNSAT answer. Put $W=\max_v|B_v|$ and $K=\max_vk_v$.

\begin{theorem}[Tree of signed factor cores; restricted exact algorithm]
For every correctly encoded source $H$ and \emph{supplied, validated}
tree certificate as above, there is one uniform exact algorithm for
both SAT and UNSAT, including arbitrary signs, clause-wise interval
demands and repeated use of the same source name in different blocks,
with
\[
T(H,T)\le \operatorname{poly}(N)
\left(1+\sum_{v\in V(T)}2^{|B_v|+k_v}\right),\qquad
S(H,T)\le \operatorname{poly}(N)+
O\!\left(\sum_{v\ne r}2^{|S_{v,p(v)}|}\right)
\]
bits in a fixed bit-machine implementation, where $r$ is the chosen root.
Consequently supplied certificates with $W+K=O(\log N)$ are
deterministically polynomial-time decidable, even when internal
block primal treewidth is unbounded.
\end{theorem}
\begin{proof}
Charge source parsing, source-name dictionaries, tree construction,
and checking that each $I_x$ is connected by a traversal of $T[I_x]$.
For every block $v$ and every assignment $b$ to the exposed original
names $B_v$, assign each of its $k_v$ higher-occurrence internal names
in all $2^{k_v}$ ways. Substitute the assigned signed literals, decrementing
both bounds of each constraint by its number of fixed TRUE literals.
Every remaining source variable now occurs in at most two constraints.
The exact signed interval $(g,f)$-factor reduction of
INTERVAL-FACTOR134, including the mandatory degree-one choice vertex
for opposite-sign occurrences, decides the residual in polynomial time.
OR across the assigned higher-occurrence names; call the result
$L_v(b)$. It is \emph{exactly} the feasibility of the original block
conditioned on $b$, with source-name consistency retained. Graph-factor
feasibility is a classical polynomial-time theorem, not a free primitive.

Root the supplied tree. For a node $v$ with parent $p$ define a Boolean
message $M_v$ indexed by assignments $a$ to $S_{v,p}$:
\[
M_v(a)=\bigvee_{\substack{b:B_v\to\{0,1\}\\
 b|_{S_{v,p}}=a}}
\left(L_v(b)\wedge\bigwedge_{u\text{ child of }v}
M_u(b|_{S_{v,u}})\right).
\]
For the root the parent separator is empty. A global satisfying source
assignment restricts to a feasible local $b$ at every block and hence
to a compatible message trajectory. Conversely a TRUE message
supplies local satisfying extensions; running intersection guarantees
that any name shared by two descendant blocks also appears throughout
their connecting tree path, so all chosen local values glue to one
global source assignment. This proves SAT correctness and, by
contraposition, UNSAT correctness without an UNSAT witness oracle.

Each candidate $b$ is enumerated at most once per block, local
residual factor feasibility costs at most $2^{k_v}\operatorname{poly}(N)$,
and child message lookups cost $\deg(v)\operatorname{poly}(N)$,
giving the stated charged bound. One completed message per nonroot
block has at most $2^{|S_{v,p(v)}|}$ bits. A polynomial-size graph-factor
algorithm's workspace is reused, along with parsers, pins, counters,
tree indices, failed branch state and buffers. The named representation
remains in live storage where needed. These statements are loose
polynomial bit-machine upper bounds, \emph{not} FLMV-small exponents.
\end{proof}

\paragraph{Structural reach, not hardness.}
One block with occurrence-two variables can encode the edge-variable
constraints of an arbitrary cubic graph $G$: every graph vertex owns
a three-literal interval clause over incident edge variables, so each
source variable occurs in exactly two clauses. For a cubic family of
unbounded treewidth, the primal graph is the line graph of $G$ and also
has unbounded treewidth. Factor feasibility remains polynomial.
The ordinary all-positive OR version is even satisfied by all-TRUE.
Thus large graph width need not force this algorithm or any other
decider to spend exponential time. The extension above adds a
tree-shaped interface around factor-solvable cores. It is a
specialization of standard tree-structured CSP message passing,
not a novel general-CSP complexity classification.

\section{Falsification and a reproducible finite test}
The independent Python reproducer generates signed interval constraints
on a four-block star with a global shared name, three spoke names and
private interior variables. Some sources have one private source
variable with three occurrences; the code explicitly branches on it
before enumerating feasible graph-factor edge subsets. Its independent
reference checker exhausts source assignments without using the
factor or tree-message code. It also tests a shared boundary forced
TRUE in one block and FALSE in another.

For three blocks containing $x\ne y$, $y\ne z$, $z\ne x$, every
individual block is feasible but their conjunction is UNSAT. No
three-vertex tree meets running intersection for their occurrence
sets. The validator \emph{rejects} the proposed path certificate;
it does not falsely apply the theorem. The global instance itself
is an easy affine/XOR control.

\textbf{COMPUTATIONAL:} seed 202610250230, 72 randomized plus two
fixed valid both-label cases: 74 valid sources, 58 SAT and 16 UNSAT,
independent source enumeration versus tree-factor decisions with
zero discrepancies; 2,024 boundary assignments, 2,155 factor calls,
13,169 factor-edge subsets inspected, 153,261 source truth rows
tested, and one invalid running-intersection certificate rejected.
Sixteen random cases included the higher-degree interior branch.
The finite checker uses brute edge-factor feasibility in place of
a production polynomial $(g,f)$-factor subroutine; this verifies
mapping correctness, \emph{not} the asymptotic bound empirically.
Python counters are not FLMV-machine primitive-bit counts.

\section{Typed state/resource ledger and hostile controls}
At state 0, the uniform controller has only encoded records, not
known clauses, a correct tree certificate or a label. Source bit
reads, signs, headers, names and paid source validation precede
all further use. At state 1 it may choose SCC, Horn, Gaussian
affine elimination, direct witness recognition, a tested factor-tree
certificate or a generic exact fallback. No finite strategy menu
is declared comprehensive of every possible decider. At state 2
the factor-tree method enumerates exact boundary values and any
high-degree branch values; it must pay all failed branches, graph
construction, choice-vertex links, array writes (even unchanged
writes), cache access, tree traversals and retained messages.
At state 3, a nonempty root message returns SAT, otherwise UNSAT;
a one-bit decider need not emit all witnesses or all local relations.

Keep distinct primitive bit work; read-only source reads; workspace
reads and writes; temporal complete-state flips; peak live writable
bits; cumulative live-bit-time; and program/representation size.
The theorem's space expression includes message tables and polynomial
factor-workspace but does not claim a small exponent. Actual flips
and bit-time were \emph{not} measured. Cached factor results or
alternative representations can improve work. On affine,
Horn, padded 2-SAT, disconnected, large-width-all-positive and
source-witness-easy families, direct recognition may win. The
tree algorithm is an \emph{upper bound} for its certified class,
never a compulsory all-decider lower bound.

\section{Gap-to-evidence acceptance and five-role reconciliation}
\begin{description}
\item[FILLED restricted:] Existing source-degree-two interval factor
lemma plus a validated source-name running-intersection tree yields
the exact message theorem above, with independent small finite checks.
The existence of a supplied certificate is part of the promise.
\item[PARTIAL/RESTRICTED:] High-degree interior branching adds the
independent parameter $k_v$, allowing arbitrary but charged switching
inside blocks. Recognition of \emph{an existing good tree certificate}
is not shown cheap if the certificate is not supplied.
\item[FALSIFIED/RETIRED:] Independent satisfiability of blocks suffices;
a cyclic three-block inequality control refutes that status-only
gluing. Width, candidate message volume, and factor-call counts as
unavoidable computation for every solver remain invalid inferences.
\item[OPEN, ARC-EDGE-90:] A source-coherent transformation which,
for every ordinary SAT input and every relevant uniform base decider
pair, realizes an all-input favorable joint time-space successor,
or a cofinal finite route, in the exact FLMV model. A cheap
normalizer into $W+K=O(\log N)$ is precisely the missing nontrivial
step; finite tests and a supplied decomposition cannot establish it.
\end{description}

\textbf{CORPUS REVIEWED.} Live repository tree (5,197 entries,
untruncated), root main.tex, main13.tex, AGENTS.md, CODEX.md,
the repository pnp-research skill, current task inventory,
original LSPV-1 INTERVAL-FACTOR134 and its handoff,
LSPV-2 IDENTITY-SWEEP129/OVERLAP-CUT128,
LSPV-3 TYPED-DEGREE-BRIDGE98,
LSPV-4 NAME-ONLY-MULTIPASS93,
LSPV-5 NAMED-FRONTIER-COMPOSITION132/STRICT-SUBJOIN-STATUS131,
and the relevant MAIN108 weighted-transcript original.
The original ARC Google Doc
\texttt{1QirzuNPCG-CrE3EX6WweqsPdDbtoblCDEE4YI0HGyC0}
was reread in 221 paragraphs. The prior pinned snapshot's
3,532-file hash/keyword census is inherited, not a fresh
full mathematical audit. Remaining unread queue includes
main2--main12 and main14 bodies, many task bodies and root
revisions, Drive LOG/PROVED\_EXACT/BOOL/SIL pagination and
the unrecovered executable original HEREDITARY-SCOPED10.
No absence is inferred from a bounded search.

\textbf{DUPLICATION / CONTRADICTION AUDIT.}
The general tree-message principle and $(g,f)$-factor solver
are \textbf{KNOWN THEOREM} components; this report gives a
source-typed and fully quantified combination, not priority over
the broader CSP literature. Role 2's exact union count
$2^{|S_i\cup S_j|}$ covers eager path pins, not necessary
factor-tree decision work. Role 3's typed equality contraction
cautions that aliases and degree counts depend on encoding.
Role 4's name-only multi-pass lower bound applies to a
sequential-input channel, not the FLMV random-access model.
Role 5's measured named-frontier repair is semantically sound
but can generate children larger than their source, so it is
not yet an ARC power shrink. None conflicts with the theorem.

\textbf{WHAT CHANGED RELATIVE TO ROOT MAIN.}
Add the signed interval-factor \emph{tree-of-cores} upper bound
and the explicit interface-plus-high-degree source parameter
$W+K$, while preserving the previous graph-factor and
named-frontier results as prerequisites. Do not promote a new
SAT lower bound, a source normalizer, or a favorable pair.

\textbf{ADOPTED:} the original degree-two factor equivalence,
source-named compatibility and FLMV exact endpoint.
\textbf{REJECTED:} treating a non-tree block incidence cycle as
a valid join certificate, treating a table or certificate
search as free, and treating a large table as mandatory work.
\textbf{RETIRED:} the historical universal cosine landing.
The next indispensable original remains the full
HEREDITARY-SCOPED10 source, followed by an all-input normalizer
test on Role 5's minimal UNSAT subjoin collisions and Role 3's
connected degree-three equality gadgets.

\section{Independent literature and transfer limitations}
\begin{thebibliography}{8}\small
\bibitem{flmv} L. Fortnow, R. Lipton, D. van Melkebeek, A. Viglas,
Time-Space Lower Bounds for Satisfiability,
\emph{J. ACM} 52(6):835--865 (2005),
\url{https://doi.org/10.1145/1101821.1101822}.
Corollary 5.8 is in the author's PDF printed page 18.
\bibitem{factor} H. N. Gabow,
Data Structures for Weighted Matching and Extensions to
$b$-Matching and $f$-Factors, \emph{ACM Trans. Algorithms}
14(3) (2018), \url{https://doi.org/10.1145/3183369}.
\bibitem{csp} G. Gottlob, G. Greco and F. Scarcello,
Tree Projections and Constraint Optimization Problems:
Fixed-Parameter Tractability and Parallel Algorithms (2017).
\url{https://arxiv.org/abs/1711.05216}.
Tree-like joining of efficiently computed subproblem views
is established prior art; detecting suitable decompositions
may itself be difficult.
\bibitem{arc} Project LSPV-1, INTERVAL-FACTOR134,
\path{research/agents/lspv-1/2026-10-25_0058_LSPV-1_INTERVAL_FACTOR134.tex}.
\bibitem{peer} Project LSPV-5, NAMED-FRONTIER-COMPOSITION132,
\path{research/agents/lspv-5/2026-10-25_0150_LSPV-5_NAMED_FRONTIER_COMPOSITION132.tex}.
\end{thebibliography}
\end{document}
% END_OF_VERIFIED_REPORT::LSPV-1::TREE-FACTOR-JOIN135
